%===================================================================================================
\section{Design Goals and Principles} %+++++++++++++++++++++++++++++++++++++++++++++++++++++++++++++
\label{sctn:design-goals-and-principles} %++++++++++++++++++++++++++++++++++++++++++++++++++++++++++
%===================================================================================================


\improve


\begin{justify}
The design of Assertify is guided by four principles that target novice-relevant
shortcomings of baseline output: clarity, contextualisation, usability, and
flexibility. These principles mirror the objectives stated for the library in
the current thesis draft, but are restated here in a form that directly
motivates concrete technical decisions.
\end{justify}


%===============================================================================
\paragraph{Clarity through a stable report schema.}%
\label{par:clarity-through-a-stable-report-schema} The primary goal is
to present failure information in a structured and easily readable form, such
that learners can quickly identify (i) what was evaluated, (ii) what was
expected, and (iii) what was actually obtained. In contrast to baseline failures
that interleave this information with low-level noise, Assertify prioritises a
compact summary view centred on an explicit expected/actual comparison.


%===============================================================================
\paragraph{Context where it supports debugging.}%
\label{par:context-where-it-supports-debugging} Some failures cannot be
understood solely from a final mismatch, especially when a failure is preceded
by intermediate evaluations or when expressions are syntactically complex.
Assertify therefore supports adding context in two complementary ways: (i)
simplifying or reducing expressions to remove irrelevant technical detail while
preserving the structure that matters for interpretation, and (ii) optionally
attaching a \emph{history} of previously evaluated expressions in scenarios
where a sequence of operations precedes a failure (e.g., stateful data structure
tests).


%===============================================================================
\paragraph{Usability and minimal intrusion.}%
\label{par:usability-and-minimal-intrusion} A central design choice is to
integrate into existing workflows without requiring a new execution environment
or substantial changes to test code. The library is therefore implemented as a
wrapper around established frameworks (e.g., MSTest assertions and FsCheck
properties) and aims to preserve familiar testing idioms while improving the
resulting feedback.


%===============================================================================
\paragraph{Flexibility and configuration.}%
\label{par:flexibility-and-configuration} Students, tutors, and developers can
have different information needs. Assertify is therefore designed to support
configurable verbosity, for example by enabling or disabling optional report
components such as reduction steps or history information. This allows the same
testing infrastructure to serve both novice-oriented feedback (minimal,
high-signal summaries) and more detailed debugging scenarios (extended context).


%===================================================================================================
% EOF ++++++++++++++++++++++++++++++++++++++++++++++++++++++++++++++++++++++++++++++++++++++++++++++
%===================================================================================================